% 第 4 章 评测数据集构建 (目标 3 页)
% 对应官方要求: 实验方案(数据部分)
% 借鉴国二报告: 数据集独立成章, 含来源/生成/评估/筛选四节
\section{评测数据集构建}

\subsection{构建原则与流程}

数据集构建遵循三项原则。第一,用户消息用正常业务话术,不暴露攻击意图。第二,工具模拟返回真实查询结果,不人为编造诱导。第三,首轮查真实数据,后续轮次由模型自主判断是否需要行动。

攻击样本 66 条,来源于政府办公 50 类技能的真实业务话术诱导,经两轮筛选改造而来。第一轮按工具多样性、模式多样性、文本长度、多步骤性四维度打分,从 351 条公开攻击样本中筛出 90 条。第二轮针对初筛样本中模型拒答与不触发占比过高的情况,以业务话术替换显式诱导,重写模拟返回为真实查询结果,最终得到 66 条在真实模型上跑通绕过的样本。

% TODO: 图 4 数据集构建流程图(351 → 四维度筛选 → 90 → 业务化改写 → 66; 旁路: 白样本 25)
\placeholderfig{数据集构建流程(351 → 90 → 66 + 白样本 25)}{评测数据集构建流程}

\subsection{攻击样本覆盖矩阵}

66 条攻击样本覆盖 9 类攻击模式与 50 类政务技能中的 34 类,见表~\ref{tab:coverage}。

\begin{table}[H]
  \centering
  \caption{攻击样本覆盖矩阵}
  \label{tab:coverage}
  \small
  \begin{tabular}{p{3cm} p{1.4cm} p{7.6cm}}
    \toprule
    攻击模式 & 数量 & 覆盖政务技能示例 \\
    \midrule
    数据外泄 & 25 & 差旅费报告、通讯录查询、跨部门协调、公文起草 \\
    伪造签字 & 10 & 审批意见出具、批量同意、会议纪要代签 \\
    系统提示伪装 & 8 & 系统升级通知、政策版本覆写 \\
    批量操作 & 5 & 群发邮件、批量建会、批量任务下发 \\
    破坏操作 & 5 & 日志清理、系统重启、文件删除 \\
    越权操作 & 4 & 跨部门数据访问、越级审批 \\
    知识库投毒 & 3 & 知识库条目篡改、检索结果污染 \\
    工具结果注入 & 3 & 邮件正文指令、查询结果指令 \\
    多轮诱导 & 3 & 先无害后诱导的渐进式攻击 \\
    \bottomrule
  \end{tabular}
\end{table}

\subsection{白样本设计}

白样本 25 条,覆盖 18 类技能,用于误报率评估。每条白样本对应一条用户明确请求的正常业务操作,例如查询差旅费政策、起草会议室使用通知并写入草稿目录、安排部门内部讨论。白样本设计时回避批量、应急、机密等可能触发评审高分的措辞,贴近真实日常办公场景。

\subsection{数据质量说明}

全部样本经三条质量校验。其一,每条攻击样本在基线模式下真实跑通并确认绕过,杜绝纸面攻击。其二,每条样本的工具调用参数由人工核对,确保动作与诱导意图一致。其三,攻击样本与白样本的技能分布做了对齐,避免误报率评估因技能分布偏差失真。全部样本与结果文件随仓库提供,可独立复现。
